% GENERATED FILE. Edit canonical lessons, then run npm run generate:books.
% canonical-source-hash: fnv1a64:567badd5a7770ec9
% canonical-lessons: ZH-C133-dasao, ZH-W133-sao-sweep, ZH-C133-jiawu, ZH-C133-shoushi, ZH-C133-ca, ZH-C133-die

\chapter{Housework}
\label{ch:zh-housework}
\hlchaptermodality{133}

\begin{chapteropening}
\textbf{By the end of this chapter:} \emph{I can say I clean, tidy up, wipe and throw things away.}

Complete the last lesson of chapter 133: I can say I clean, tidy up, wipe and throw things away.
\end{chapteropening}

\section[\texorpdfstring{\zh{打扫} (dǎsǎo)}{打扫 (dǎsǎo)}]{\zh{打扫} (dǎsǎo) — to clean (a room); to sweep}

\label{lesson:ZH-C133-dasao}



\begin{quote}
Before the new one: say the Mandarin for in all, then the Mandarin for to count.
\end{quote}

\subsection*{You'll want to know: \zh{打扫}}
\textbf{\zh{打扫}} — \emph{dǎsǎo} — \textquotedblleft{}to clean (a room); to sweep\textquotedblright{}.

\begin{quote}
\textbf{\zh{今天我们一起打扫}} \emph{jīntiān wǒmen yìqǐ dǎsǎo} — Today we're cleaning up together.
\end{quote}

The new character, in pieces: \zh{扫} — the hand piece on the left, as in \zh{打}, and on the right three strokes: a bar that turns down, and two bars inside it. The next lesson puts it in your hand.

Two third tones: dǎsǎo is said dásǎo, the first rising.

\begin{cousinweb}
\textbf{\zh{打扫}} (\emph{dǎsǎo}) is built from \zh{打} 'to strike', here 'to do', and \emph{sǎo} 'to sweep', both inherited from Middle Chinese, \emph{taengX} and \emph{sawX} (Guangyun, in Baxter's transcription). The regular sound change from \emph{taengX} would not give the everyday \emph{dǎ}; why is not settled.
\end{cousinweb}

\begin{grammarlens}[title={What you've built}]
One more word for everyday actions.
\end{grammarlens}

\subsection*{Your turn}
\begin{itemize}
\raggedright
  \item \emph{Say it:} \emph{dǎsǎo}
  \item \emph{Say it:} \emph{dǎsǎo}, once more
  \item \emph{Say it:} \emph{jīntiān wǒmen yìqǐ dǎsǎo} — Today we're cleaning up together.
  \item \emph{Recall:} say the Mandarin for in all, then the Mandarin for to count, then say \emph{dǎsǎo} again
\end{itemize}

\begin{tcolorbox}[breakable,colback=teal!4,colframe=teal!35!black,title={Before you move on}]
What does \zh{打扫} mean? (\textquotedblleft{}To clean (a room); to sweep\textquotedblright{}.) Say it once more.
\end{tcolorbox}
\section[\texorpdfstring{\zh{扫} (sǎo)}{扫 (sǎo)}]{\zh{扫} — the hand piece, then three bars}

\label{lesson:ZH-W133-sao-sweep}



\begin{quote}
The word first, then the character inside it.

\begin{itemize}
\raggedright
  \item \emph{Recall:} say \emph{dǎsǎo} — \textbf{R1}, one lesson back, before it decays
\end{itemize}

\begin{quote}
\zh{打} \zh{手}
\end{quote}

\emph{Read it:} each one aloud, with nothing said first — \emph{dǎ}, \emph{shǒu}
\end{quote}

\subsection*{Script — the hand piece, then three bars}
\hlblockfigure{\detokenize{figures/ZH-W133-sao-sweep-filmstrip.pdf}}{How \zh{扫} is written, stroke by stroke}

\begin{quote}
\zh{扫}
\end{quote}

\textbf{Two pieces, left before right.}

\textbf{The hand piece, three strokes}, as in \zh{打}: a horizontal, a vertical that hooks to the left at its foot, then a stroke that rises to the right across it.

\textbf{The right piece, three strokes.} A horizontal that turns down at its right end, a horizontal across the middle, then a horizontal along the foot, which runs out past the right side.

3 + 3 = six.

\subsection*{Your turn}
\begin{itemize}
\raggedright
  \item \emph{Write it:} \zh{扫}, all six, with the model in front of you
  \item \emph{Cover:} write it again from memory, and count the pieces, not the strokes
\end{itemize}

\begin{tcolorbox}[breakable,colback=teal!4,colframe=teal!35!black,title={Before you move on}]
Which piece of \zh{扫} comes from \zh{打}? (\textbf{The hand piece, on the left}.) How many strokes on the right? (\textbf{Three}.) Which bar runs out past the side? (\textbf{The bottom one}.)

Source: \href{https://raw.githubusercontent.com/chanind/hanzi-writer-data/68d10a4b21150cae5e1ebbd223eed289cf32d90c/data/\%E6\%89\%AB.json}{Hanzi Writer Data \zh{扫}, pinned stroke paths and medians}.
\end{tcolorbox}
\section[jiāwù]{jiāwù — housework}

\label{lesson:ZH-C133-jiawu}



\begin{quote}
Before the new one: say the Mandarin for to count, then the Mandarin for to clean (a room).
\end{quote}

\subsection*{You'll want to know: jiāwù}
\emph{jiāwù} — \textquotedblleft{}housework\textquotedblright{}.

\begin{quote}
\emph{wǒ bù xǐhuan zuò jiāwù} — I don't like doing housework.
\end{quote}

Its characters come later; for now it is a word for the ear.

\begin{cousinweb}
\emph{jiāwù} is built from \zh{家} 'home' and \emph{wù} 'affairs', both inherited from Middle Chinese, \emph{kae} and \emph{mjuH} (Guangyun, in Baxter's transcription).
\end{cousinweb}

\begin{grammarlens}[title={What you've built}]
One more word for everyday actions.
\end{grammarlens}

\subsection*{Your turn}
\begin{itemize}
\raggedright
  \item \emph{Say it:} \emph{jiāwù}
  \item \emph{Say it:} \emph{jiāwù}, once more
  \item \emph{Say it:} \emph{wǒ bù xǐhuan zuò jiāwù} — I don't like doing housework.
  \item \emph{Recall:} say the Mandarin for to count, then the Mandarin for to clean (a room), then say \emph{jiāwù} again
\end{itemize}

\begin{tcolorbox}[breakable,colback=teal!4,colframe=teal!35!black,title={Before you move on}]
What does jiāwù mean? (\textquotedblleft{}Housework\textquotedblright{}.) Say it once more.
\end{tcolorbox}
\section[shōushi]{shōushi — to tidy up}

\label{lesson:ZH-C133-shoushi}



\begin{quote}
Before the new one: say the Mandarin for to clean (a room), then the Mandarin for housework.
\end{quote}

\subsection*{You'll want to know: shōushi}
\emph{shōushi} — \textquotedblleft{}to tidy up\textquotedblright{}.

\begin{quote}
\emph{jīntiān wǒ shōushi fángjiān} — Today I'm tidying my room.
\end{quote}

Its characters come later; for now it is a word for the ear.

\begin{cousinweb}
\emph{shōushi} is built from \emph{shōu} 'to gather' and \emph{shi}, both inherited from Middle Chinese, \emph{syuw} and \emph{dzyip} (Guangyun, in Baxter's transcription); the dictionary glosses the second with this very word, and with 'to pick up'.
\end{cousinweb}

\begin{grammarlens}[title={What you've built}]
One more word for everyday actions.
\end{grammarlens}

\subsection*{Your turn}
\begin{itemize}
\raggedright
  \item \emph{Say it:} \emph{shōushi}
  \item \emph{Say it:} \emph{shōushi}, once more
  \item \emph{Say it:} \emph{jīntiān wǒ shōushi fángjiān} — Today I'm tidying my room.
  \item \emph{Recall:} say the Mandarin for to clean (a room), then the Mandarin for housework, then say \emph{shōushi} again
\end{itemize}

\begin{tcolorbox}[breakable,colback=teal!4,colframe=teal!35!black,title={Before you move on}]
What does shōushi mean? (\textquotedblleft{}To tidy up\textquotedblright{}.) Say it once more.
\end{tcolorbox}
\section[cā]{cā — to wipe}

\label{lesson:ZH-C133-ca}



\begin{quote}
Before the new one: say the Mandarin for housework, then the Mandarin for to tidy up.
\end{quote}

\subsection*{You'll want to know: cā}
\emph{cā} — \textquotedblleft{}to wipe\textquotedblright{}.

\begin{quote}
\emph{wǒ cā zhuōzi nǐ cā yǐzi} — I'll wipe the table; you wipe the chairs.
\end{quote}

Its characters come later; for now it is a word for the ear.

\begin{cousinweb}
\emph{cā} No dictionary we use records an older form, so its history is not settled.
\end{cousinweb}

\begin{grammarlens}[title={What you've built}]
One more word for everyday actions.
\end{grammarlens}

\subsection*{Your turn}
\begin{itemize}
\raggedright
  \item \emph{Say it:} \emph{cā}
  \item \emph{Say it:} \emph{cā}, once more
  \item \emph{Say it:} \emph{wǒ cā zhuōzi nǐ cā yǐzi} — I'll wipe the table; you wipe the chairs.
  \item \emph{Recall:} say the Mandarin for housework, then the Mandarin for to tidy up, then say \emph{cā} again
\end{itemize}

\begin{tcolorbox}[breakable,colback=teal!4,colframe=teal!35!black,title={Before you move on}]
What does cā mean? (\textquotedblleft{}To wipe\textquotedblright{}.) Say it once more.
\end{tcolorbox}
\section[dié]{dié — to fold}

\label{lesson:ZH-C133-die}



\begin{quote}
Before the new one: say the Mandarin for to tidy up, then the Mandarin for to wipe.
\end{quote}

\subsection*{You'll want to know: dié}
\emph{dié} — \textquotedblleft{}to fold\textquotedblright{}.

\begin{quote}
\emph{wǒ dié yīfu nǐ shōushi fángjiān} — I'll fold the clothes; you tidy the room.
\end{quote}

Its characters come later; for now it is a word for the ear.

\begin{cousinweb}
\emph{dié} is inherited from Middle Chinese \emph{dep} (Guangyun, in Baxter's transcription), glossed 'layered' and 'to pile up'; folding, a cloth laid in layers, is the likely step, though it is not given there.
\end{cousinweb}

\begin{grammarlens}[title={What you've built}]
One more word for everyday actions.
\end{grammarlens}

\subsection*{Your turn}
\begin{itemize}
\raggedright
  \item \emph{Say it:} \emph{dié}
  \item \emph{Say it:} \emph{dié}, once more
  \item \emph{Say it:} \emph{wǒ dié yīfu nǐ shōushi fángjiān} — I'll fold the clothes; you tidy the room.
  \item \emph{Recall:} say the Mandarin for to tidy up, then the Mandarin for to wipe, then say \emph{dié} again
\end{itemize}

\begin{tcolorbox}[breakable,colback=teal!4,colframe=teal!35!black,title={Before you move on}]
What does dié mean? (\textquotedblleft{}To fold\textquotedblright{}.) Say it once more.
\end{tcolorbox}
